\documentclass[12pt,a4paper]{article}
% Drake_preamble.tex - LaTeX Preamble Template
% 作者: 謝慕揚, MD, PhD, FESC
% 
% 使用說明:
% 1. 表格請使用 longtable，不要有垂直分隔線 |
% 2. 表格內換行使用 \newline，不要使用 \\
% 3. 藥物名稱、檢驗、檢查請維持英文
% 4. Reference 格式請使用 \href 加上驗證過的 DOI

% Including essential packages for Chinese typesetting
\usepackage{xeCJK}
\usepackage{fontspec}
% Configuring Chinese font with Noto Serif CJK TC, fallback to SimSun
\IfFontExistsTF{Noto Serif CJK TC}{
  \setCJKmainfont{Noto Serif CJK TC}
  \setCJKsansfont{Noto Serif CJK TC}
  \setCJKmonofont{Noto Serif CJK TC}
}{\setCJKmainfont{SimSun}
  \setCJKsansfont{SimSun}
  \setCJKmonofont{SimSun}
}

% Including other necessary packages
\usepackage{geometry}
\usepackage{titlesec}
\usepackage{enumitem}
\usepackage{xcolor}
\usepackage{colortbl}
\usepackage{tcolorbox}
\usepackage{booktabs}
\usepackage{longtable}
\usepackage{array}
\usepackage{graphicx}
\usepackage{tikz}
\usepackage{fancyhdr}
\usepackage{multicol}
\usepackage{datetime2}
\usepackage{hyperref}
\usepackage{mdframed}
\usepackage{amsmath}
\usepackage{amssymb}
\usepackage{multirow}

\hypersetup{
    colorlinks=true,
    linkcolor=emergency_blue,
    citecolor=emergency_blue,
    urlcolor=emergency_blue,
    pdfborder={0 0 0}
}

% Defining page geometry
\geometry{
  top=2.5cm,
  bottom=2.5cm,
  left=2cm,
  right=2cm,
  headheight=25pt
}

% Adjusting topmargin to compensate for increased headheight
\addtolength{\topmargin}{-10pt}

% Defining colors
\definecolor{emergency_red}{RGB}{186,24,27}
\definecolor{emergency_blue}{RGB}{0,114,188}
\definecolor{emergency_gray}{RGB}{120,120,120}
\definecolor{highlight_yellow}{RGB}{255,250,205}

% Setting title formats
\titleformat{\section}[block]{\Large\bfseries\color{emergency_red}}{\thesection}{10pt}{}
\titleformat{\subsection}[block]{\large\bfseries\color{emergency_blue}}{\thesubsection}{8pt}{}
\titleformat{\subsubsection}[block]{\normalsize\bfseries\color{emergency_gray}}{\thesubsubsection}{6pt}{}

% Defining custom environment for important notes
\newmdenv[
    linecolor=emergency_red,
    backgroundcolor=highlight_yellow,
    linewidth=2pt,
    topline=true,
    bottomline=true,
    leftline=true,
    rightline=true,
    innertopmargin=10pt,
    innerbottommargin=10pt,
    innerleftmargin=10pt,
    innerrightmargin=10pt
]{importantbox}

% Defining note command with tcolorbox
\newcommand{\note}[1]{
    \par
    \noindent
    \begin{tcolorbox}[
        colback=emergency_blue!5!white,
        colframe=emergency_blue,
        title=\textbf{重點提醒},
        fonttitle=\bfseries,
        boxrule=0.5pt,
        arc=3pt,
        left=6pt,
        right=6pt,
        top=6pt,
        bottom=6pt
    ]
    #1
    \end{tcolorbox}
}

% Key findings box
\newcommand{\keyfinding}[1]{
    \par
    \noindent
    \begin{tcolorbox}[
        colback=yellow!10!white,
        colframe=orange!75!black,
        title=\textbf{核心發現摘要},
        fonttitle=\bfseries,
        boxrule=0.5pt,
        arc=3pt
    ]
    #1
    \end{tcolorbox}
}

% Defining custom date format for ISO 8601
\DTMnewdatestyle{iso}{
  \renewcommand{\DTMdisplaydate}[4]{\number##1-\number##2-\number##3}
  \renewcommand{\DTMDisplaydate}[4]{\number##1-\number##2-\number##3}
}
\DTMsetdatestyle{iso}
\newcommand{\mydate}{\DTMtoday}


\pagestyle{fancy}
\fancyhf{}
\fancyhead[L]{\textcolor{emergency_red}{EuroIntervention 教學講義}}
\fancyhead[R]{\textcolor{emergency_blue}{2026年1月號精選}}
\fancyfoot[C]{\textcolor{emergency_gray}{\thepage}}
\renewcommand{\headrulewidth}{1pt}
\renewcommand{\headrule}{\hbox to\headwidth{\color{emergency_red}\leaders\hrule height \headrulewidth\hfill}}

\begin{document}

\begin{center}
{\LARGE\bfseries\textcolor{emergency_red}{EuroIntervention}}\\[0.5cm]
{\Large\textbf{January 2026 Issue Snapshot (Volume 22, Issue 1)}}\\[0.3cm]
{\large 2026年1月號期刊精選}\\[0.5cm]
{\normalsize \href{https://eurointervention.pcronline.com/issue/volume-22/number-1}{EuroIntervention 2026;22(1)}}\\[0.3cm]
{\normalsize 整理：謝慕揚 MD, PhD, FESC \quad 日期：\mydate}
\end{center}

\vspace{0.5cm}

\begin{tcolorbox}[colback=yellow!10!white,colframe=orange!75!black,title=\textbf{本期核心重點}]
\begin{enumerate}
    \item \textbf{Expert Review：}Endomyocardial biopsy (EMB)的最佳化臨床應用——從患者選擇到技術步驟
    \item \textbf{新技術：}IMR$_{CT}$——基於CCTA的非侵入性冠狀動脈微血管疾病評估
    \item \textbf{STEMI預後：}OCT研究顯示non-ruptured thin-cap fibroatheroma（而非non-culprit plaque rupture）是長期不良預後的獨立預測因子
    \item \textbf{QFR再現性：}FAVOR III Europe的REPEAT-QFR substudy揭示in-procedure與core lab QFR分析的一致性僅為中等
\end{enumerate}
\end{tcolorbox}

\tableofcontents
\newpage

%============================================================
\section{Expert Review：Endomyocardial Biopsy最佳化}
%============================================================

\subsection{EMB的臨床適應症}

Endomyocardial biopsy (EMB)用於診斷多種心臟疾病：

\begin{itemize}
    \item \textbf{心臟移植排斥監測}
    \item \textbf{心肌炎}（包括giant cell myocarditis、eosinophilic myocarditis）
    \item \textbf{浸潤性心肌病}（cardiac amyloidosis、hemochromatosis）
    \item \textbf{心臟腫瘤}
    \item \textbf{Arrhythmogenic cardiomyopathy}
    \item \textbf{不明原因心衰}
\end{itemize}

\subsection{重要技術考量}

\begin{longtable}{p{4cm}p{10cm}}
\toprule
\textbf{項目} & \textbf{內容} \\
\midrule
\endfirsthead
\toprule
\textbf{項目} & \textbf{內容} \\
\midrule
\endhead
\bottomrule
\endfoot

部位選擇 & \textbf{RV EMB}：較常用、較安全\newline
\textbf{LV EMB}：特定情況需要（如LV-predominant disease） \\
\midrule
影像導引 & Fluoroscopy、TTE、ICE可輔助定位 \\
\midrule
取樣數量 & 通常4-6個samples以提高診斷率 \\
\midrule
主要併發症 & 心臟穿孔（$<$1\%）、心律不整、三尖瓣損傷、血管穿刺併發症 \\
\bottomrule
\end{longtable}

\note{
\textbf{未來展望：}EMB結合advanced molecular diagnostics（如gene expression profiling、next-generation sequencing）可提供更精準的診斷與治療指引。
}

%============================================================
\section{新技術：IMR$_{CT}$——非侵入性微血管功能評估}
%============================================================

\subsection{背景}

\begin{itemize}
    \item \textbf{Coronary microvascular disease (CMD)}是重要但常被忽略的疾病
    \item 傳統評估需侵入性方法（如IMR、CFR測量）
    \item 需要可靠的\textbf{非侵入性}評估方法
\end{itemize}

\subsection{IMR$_{CT}$方法學}

\begin{longtable}{p{4cm}p{10cm}}
\toprule
\textbf{項目} & \textbf{內容} \\
\midrule
\endfirsthead
\bottomrule
\endfoot

技術基礎 & 基於coronary CT angiography (CCTA) \\
\midrule
資料整合 & Multiphase CCTA + retrospective ECG gating \\
\midrule
評估指標 & Computed tomography-based index of microvascular resistance (IMR$_{CT}$) \\
\midrule
驗證方式 & 與侵入性IMR測量比較 \\
\bottomrule
\end{longtable}

\subsection{研究結果}

\keyfinding{
IMR$_{CT}$與侵入性IMR在\textbf{vessel level}和\textbf{patient level}均呈現良好相關性，提供了一個可行的非侵入性CMD評估工具。
}

\begin{tcolorbox}[colback=emergency_blue!5!white,colframe=emergency_blue,title=\textbf{臨床意義}]
\begin{itemize}
    \item 可能成為CMD篩檢的第一線工具
    \item 減少不必要的侵入性檢查
    \item 適用於INOCA（Ischemia with No Obstructive Coronary Arteries）患者評估
\end{itemize}
\end{tcolorbox}

%============================================================
\section{STEMI研究：Non-Culprit Plaque Rupture的臨床影響}
%============================================================

\subsection{研究設計}

\begin{longtable}{p{4cm}p{10cm}}
\toprule
\textbf{項目} & \textbf{內容} \\
\midrule
\endfirsthead
\bottomrule
\endfoot

研究方法 & Optical coherence tomography (OCT) \\
\midrule
掃描範圍 & 三條主要心外膜冠狀動脈 \\
\midrule
研究族群 & STEMI患者 \\
\midrule
主要分析 & Non-culprit plaque rupture對長期預後的影響 \\
\bottomrule
\end{longtable}

\subsection{主要發現}

\begin{itemize}
    \item 有\textbf{non-culprit plaque rupture}的患者確實有較高的長期不良事件風險
    \item 但\textbf{關鍵發現}：
\end{itemize}

\keyfinding{
\textbf{Non-ruptured thin-cap fibroatheroma (TCFA)}——而非non-culprit plaque rupture——才是長期不良預後的\textbf{獨立預測因子}。
}

\subsection{臨床意義}

\begin{itemize}
    \item STEMI患者應注意non-culprit lesions的特徵
    \item TCFA的存在代表「vulnerable patient」而非僅是vulnerable plaque
    \item 可能需要更積極的藥物治療（如高強度statin、colchicine）
    \item 是否需要對TCFA進行預防性介入仍需更多研究
\end{itemize}

%============================================================
\section{方法學研究：QFR的再現性——REPEAT-QFR Substudy}
%============================================================

\subsection{研究背景}

\begin{itemize}
    \item \textbf{Quantitative Flow Ratio (QFR)}是基於血管造影的FFR替代指標
    \item \textbf{FAVOR III Europe}試驗對QFR的可靠性提出質疑
    \item 需要了解差異的來源
\end{itemize}

\subsection{REPEAT-QFR設計}

比較\textbf{in-procedure QFR}（術中即時分析）與\textbf{core lab QFR}（核心實驗室分析）的一致性。

\subsection{主要發現}

\begin{longtable}{p{5cm}p{9cm}}
\toprule
\textbf{發現} & \textbf{說明} \\
\midrule
\endfirsthead
\bottomrule
\endfoot

整體一致性 & \textcolor{red}{\textbf{中等 (Modest)}} \\
\midrule
影響變異的因素 & • Angiographic quality（影像品質）\newline
• Advanced CAD（進展性冠狀動脈疾病）\newline
• Poor in-procedure QFR analysis（術中分析品質不佳） \\
\bottomrule
\end{longtable}

\note{
\textbf{Editorial觀點（Lansky）：}此研究揭示QFR在真實世界應用時，操作者經驗與影像品質對結果有重大影響。在複雜病灶或影像品質不佳時，應謹慎解讀QFR結果，考慮使用invasive FFR確認。
}

%============================================================
\section{其他本期重點}
%============================================================

\subsection{High-Risk Plaque Interventions}

\textbf{作者：Gary S. Mintz \& Carlos Collet}

\begin{itemize}
    \item 探討如何識別與處理high-risk plaque
    \item 結合intravascular imaging的策略
    \item PCI最佳化與plaque modification技術
\end{itemize}

\subsection{STEMI in Large Coronary Arteries}

\textbf{作者：Marc Bonnet, Grégoire Rangé等}

\begin{itemize}
    \item 大血管（large coronary arteries）STEMI的處理挑戰
    \item Stent sizing考量
    \item No-reflow現象的預防與處理
    \item 長期預後分析
\end{itemize}

%============================================================
\section{2025年度審稿人致謝}
%============================================================

本期包含對EuroIntervention審稿人的年度致謝名單，表彰其對期刊與介入心臟學領域的貢獻。

\textbf{Gold Reviewers（年度10篇以上審稿）}包括：
\begin{itemize}
    \item Fernando Alfonso, Ziad Ali, Mohamed Abdel-Wahab
    \item Bernard De Bruyne, Carlos Collet, Giuseppe Tarantini
    \item Nicolas Van Mieghem, Patrick Serruys, Lorenz Räber
    \item 以及許多其他國際知名專家
\end{itemize}

%============================================================
\section{縮寫對照表}
%============================================================

\begin{longtable}{p{4cm}p{10cm}}
\toprule
\textbf{縮寫} & \textbf{全名} \\
\midrule
\endfirsthead
\bottomrule
\endfoot

EMB & Endomyocardial Biopsy (心內膜心肌切片) \\
RV & Right Ventricle (右心室) \\
LV & Left Ventricle (左心室) \\
CMD & Coronary Microvascular Disease (冠狀動脈微血管疾病) \\
IMR & Index of Microvascular Resistance (微血管阻力指數) \\
IMR$_{CT}$ & CT-based IMR (電腦斷層微血管阻力指數) \\
CCTA & Coronary CT Angiography (冠狀動脈電腦斷層血管攝影) \\
OCT & Optical Coherence Tomography (光學相干斷層掃描) \\
TCFA & Thin-Cap Fibroatheroma (薄帽纖維粥瘤) \\
QFR & Quantitative Flow Ratio (定量血流比值) \\
FFR & Fractional Flow Reserve (血流儲備分數) \\
STEMI & ST-Elevation Myocardial Infarction (ST段上升心肌梗塞) \\
INOCA & Ischemia with No Obstructive Coronary Arteries \\
CAD & Coronary Artery Disease (冠狀動脈疾病) \\
ICE & Intracardiac Echocardiography (心內超音波) \\
TTE & Transthoracic Echocardiography (經胸心臟超音波) \\
\bottomrule
\end{longtable}

%============================================================
\section{參考文獻}
%============================================================

\begin{enumerate}
    \item Dangas G, et al. Endomyocardial Biopsy: Optimising Clinical Practice. \textit{EuroIntervention}. 2026;22(1). [\href{https://eurointervention.pcronline.com/issue/volume-22/number-1}{Issue Link}]

    \item Driessen MMD, et al. IMR$_{CT}$: Non-invasive Assessment of Coronary Microvascular Resistance from CT. \textit{EuroIntervention}. 2026;22(1). [\href{https://eurointervention.pcronline.com/issue/volume-22/number-1}{Issue Link}]

    \item Kubo T, et al. Non-Culprit Plaque Rupture and Thin-Cap Fibroatheroma in STEMI: An OCT Study. \textit{EuroIntervention}. 2026;22(1). [\href{https://eurointervention.pcronline.com/issue/volume-22/number-1}{Issue Link}]

    \item Tomaniak M, et al. REPEAT-QFR: Reproducibility of Quantitative Flow Ratio in FAVOR III Europe. \textit{EuroIntervention}. 2026;22(1). [\href{https://eurointervention.pcronline.com/issue/volume-22/number-1}{Issue Link}]

    \item Mintz GS, Collet C. High-Risk Plaque Interventions: An Expert Review. \textit{EuroIntervention}. 2026;22(1). [\href{https://eurointervention.pcronline.com/issue/volume-22/number-1}{Issue Link}]

    \item Bonnet M, Rangé G, et al. STEMI in Large Coronary Arteries: Challenges and Outcomes. \textit{EuroIntervention}. 2026;22(1). [\href{https://eurointervention.pcronline.com/issue/volume-22/number-1}{Issue Link}]
\end{enumerate}

\end{document}
